% !PRINTED 72
\weberpage{72}
%<<CONT

cannot change anything; thus for \(\Delta\) only the value zero remains.

One only expresses theorem V differently when one says:

\gesperrt{VI. One obtains a vanishing determinant if the elements of one row are replaced by the corresponding elements of another row, or, shortly, if one lower or upper index is replaced by another.}

\begin{center}
{\large \S.\ 22.}

\vspace{0.5em}
{\large\gesperrt{Subdeterminants.}}
\end{center}

In every term of the expanded determinant
\[
 \sum \pm a_1^{(1)}a_2^{(2)}\cdots a_n^{(n)},
\]
whose value we now wish to denote by \(A\), every one of the numbers \(1,2,\ldots,n\) occurs once and only once as a lower index. Thus a certain complex of terms will contain the factor \(a_1^{(1)}\), another complex the factor \(a_1^{(2)}\), and so on, finally a complex the factor \(a_1^{(n)}\); every term of the determinant occurs in one and only one of these complexes.

If, therefore, we denote the first of these complexes by \(a_1^{(1)}A_1^{(1)}\), the second by \(a_1^{(2)}A_1^{(2)}\), and the last by \(a_1^{(n)}A_1^{(n)}\), then the determinant can be represented in the following way:
\begin{equation}
 A=a_1^{(1)}A_1^{(1)}+a_1^{(2)}A_1^{(2)}+\cdots+a_1^{(n)}A_1^{(n)}.
\tag{1}
\end{equation}
Instead of the lower index \(1\), we could equally well single out any other index \(\nu\), and hence put
\begin{equation}
 A=a_\nu^{(1)}A_\nu^{(1)}+a_\nu^{(2)}A_\nu^{(2)}+\cdots+a_\nu^{(n)}A_\nu^{(n)}.
\tag{2}
\end{equation}
Here the product \(a_\nu^{(\mu)}A_\nu^{(\mu)}\) denotes the complex of all terms of the determinant which contain the factor \(a_\nu^{(\mu)}\).

Since the same rules hold for upper as for lower indices, the determinant can also be written in the following way:
\begin{equation}
 A=a_1^{(\mu)}A_1^{(\mu)}+a_2^{(\mu)}A_2^{(\mu)}
 +\cdots+a_n^{(\mu)}A_n^{(\mu)},
\tag{3}
\end{equation}
where \(\mu\), too, may denote any one of the indices \(1,2,\ldots,n\).
